\documentclass[12pt,openany,oneside,a4paper,english,brazil]{abntex2}

% Pacotes fundamentais
\usepackage{lmodern}			        % Usa a fonte Latin Modern			
\usepackage[T1]{fontenc}		        % Selecao de codigos de fonte
\usepackage[utf8]{inputenc}		        % Codificacao UTF-8
\usepackage{indentfirst}		        % Indenta o primeiro paragrafo de cada secao
\usepackage{color}				        % Controle das cores
\usepackage{graphicx}			        % Inclusao de graficos
\usepackage{microtype} 			        % Melhorias de justificacao
\usepackage{amsmath, amsfonts, amssymb} % Pacotes matematicos
\usepackage{algorithm}
\usepackage{algorithmic}
\usepackage{listings}
\usepackage{booktabs}

% Pacotes de citacoes ABNT
\usepackage[brazilian,hyperpageref]{backref}
\usepackage[alf]{abntex2cite}

% Traducao e Customizacao do Ambiente de Algoritmo
\floatname{algorithm}{Algoritmo}
\renewcommand{\algorithmicrequire}{\textbf{Entrada:}}
\renewcommand{\algorithmicensure}{\textbf{Saída:}}
\renewcommand{\algorithmicfor}{\textbf{para}}
\renewcommand{\algorithmicdo}{\textbf{faça}}
\renewcommand{\algorithmicif}{\textbf{se}}
\renewcommand{\algorithmicthen}{\textbf{então}}
\renewcommand{\algorithmicelse}{\textbf{senão}}
\renewcommand{\algorithmicendif}{\textbf{fim se}}
\renewcommand{\algorithmicendfor}{\textbf{fim para}}

% Informacoes de dados para CAPA e FOLHA DE ROSTO
\titulo{GYMBEAT: SISTEMA VESTÍVEL E MÓVEL PARA RECONHECIMENTO E CONTAGEM DE REPETIÇÕES EM EXERCÍCIOS RESISTIDOS BASEADO EM SENSORES INERCIAIS E INTELIGÊNCIA ARTIFICIAL}
\autor{José Pascoal Martins}
\local{Campinas - SP}
\data{2026}
\orientador{Prof. Dr. Antonio Carlos Forster}
\instituicao{%
  PONTIFÍCIA UNIVERSIDADE CATÓLICA DE CAMPINAS - PUC-CAMPINAS
  \par
  CENTRO DE CIÊNCIAS EXATAS, AMBIENTAIS E DE TECNOLOGIAS - CEATEC
  \par
  FACULDADE DE ENGENHARIA DE COMPUTAÇÃO}
\tipotrabalho{Trabalho de Conclusão de Curso (Graduação)}
\preambulo{Trabalho de Conclusão de Curso apresentado à Faculdade de Engenharia de Computação da Pontifícia Universidade Católica de Campinas como requisito parcial para obtenção do título de Bacharel em Engenharia de Computação.}

% Customizacao da Capa (Padrao PUC-Campinas / ABNT)
\renewcommand{\imprimircapa}{%
  \begin{capa}%
    \center
    {\ABNTEXchapterfont\bfseries\normalsize \imprimirinstituicao}
    
    \vspace*{3.5cm}
    {\ABNTEXchapterfont\bfseries\Large \imprimirautor}
    
    \vspace*{3.5cm}
    {\ABNTEXchapterfont\bfseries\Large \imprimirtitulo}
    
    \vfill
    {\ABNTEXchapterfont\normalsize \imprimirlocal\\
    \imprimirdata}
    \vspace*{0.5cm}
  \end{capa}
}

% Customizacao da Folha de Rosto
\makeatletter
\renewcommand{\folhaderostocontent}{
  \begin{center}
    {\ABNTEXchapterfont\bfseries\Large \imprimirautor}
    
    \vspace*{3.5cm}
    {\ABNTEXchapterfont\bfseries\Large \imprimirtitulo}
    
    \vspace*{2.5cm}
    \hspace{.45\textwidth}
    \begin{minipage}{.5\textwidth}
      \SingleSpacing
      \imprimirpreambulo
      \vspace*{0.8cm}
      
      {\small\imprimirorientadorRotulo~{\imprimirorientador}}
    \end{minipage}%
    \vspace*{\fill}
    
    {\ABNTEXchapterfont\normalsize \imprimirlocal\\
    \imprimirdata}
    \vspace*{0.5cm}
  \end{center}
}
\makeatother

% Configuracoes de aparencia do PDF final
\definecolor{blue}{RGB}{41,5,195}
\hypersetup{
    colorlinks=true,
    linkcolor=blue,
    citecolor=blue,
    filecolor=magenta,
    urlcolor=blue
}

% Inicio do documento
\begin{document}

\selectlanguage{brazil}
\frenchspacing 

% ELEMENTOS PRE-TEXTUAIS
\imprimircapa
\imprimirfolhaderosto

% RESUMO
\setlength{\absparsep}{18pt}
\begin{resumo}
O monitoramento automático de exercícios físicos resistidos (musculação) por meio de dispositivos vestíveis (\textit{wearables}) representa uma área emergente no campo do Reconhecimento de Atividades Humanas (\textit{Human Activity Recognition} - HAR). No entanto, a maioria das soluções comerciais restringe-se a estimativas macroscópicas de gasto calórico e frequência cardíaca, carecendo de mecanismos precisos para contagem de repetições e análise biomecânica em tempo real. Este trabalho apresenta o \textbf{Gymbeat}, um sistema integrado composto por um aplicativo para \textit{smartwatch} (Apple Watch) e um aplicativo para \textit{smartphone} (iPhone) desenvolvido em React Native, destinado à coleta, processamento e contagem automática de repetições em tempo real. O sistema utiliza dados inerciais de acelerômetro e giroscópio coletados a uma frequência de 25 Hz. Para a contagem, adota-se uma abordagem híbrida fundamentada em Processamento Digital de Sinais (filtragem passa-baixa Butterworth e cálculo de magnitude invariante à rotação) combinada a um algoritmo de detecção de extremos locais e validação por Aprendizado de Máquina clássico. A metodologia de avaliação fundamenta-se no protocolo \textit{Leave-One-Subject-Out} (LOSO), avaliando métricas de Erro Médio Absoluto ($MAE$), Raiz do Erro Quadrático Médio ($RMSE$) e latência computacional. Os resultados preliminares demonstram a viabilidade técnica da arquitetura proposta, garantindo contagem confiável com latência inferior a 5 ms no dispositivo móvel.

\textbf{Palavras-chave:} Dispositivos Vestíveis. Sensores Inerciais. Reconhecimento de Atividades Humanas. Processamento Digital de Sinais. Aprendizado de Máquina. Contagem de Repetições.
\end{resumo}

% ABSTRACT
\begin{resumo}[Abstract]
\begin{otherlanguage*}{english}
Automatic tracking of resistance exercises (strength training) using wearable devices represents a key emerging area in Human Activity Recognition (HAR). However, most commercial solutions remain limited to macroscopic estimates of caloric expenditure and heart rate, lacking precise mechanisms for real-time repetition counting and biomechanical analysis. This work presents \textbf{Gymbeat}, an integrated system comprising a smartwatch application (Apple Watch) and a smartphone application (iPhone) developed in React Native, designed for real-time collection, processing, and automatic repetition counting. The system leverages inertial data from accelerometer and gyroscope sampled at 25 Hz. For repetition counting, a hybrid approach is employed, combining Digital Signal Processing (Butterworth low-pass filtering and rotation-invariant magnitude calculation) with local extrema detection and validation through classical Machine Learning. The evaluation methodology is based on the Leave-One-Subject-Out (LOSO) cross-validation protocol, assessing Mean Absolute Error ($MAE$), Root Mean Squared Error ($RMSE$), and computational latency. Preliminary results confirm the technical feasibility of the proposed architecture, delivering reliable counting with latency below 5 ms on the mobile device.

\textbf{Keywords:} Wearable Devices. Inertial Sensors. Human Activity Recognition. Digital Signal Processing. Machine Learning. Repetition Counting.
\end{otherlanguage*}
\end{resumo}

% Lista de abreviaturas e siglas
\begin{siglas}
  \item[ABNT] Associação Brasileira de Normas Técnicas
  \item[API] \textit{Application Programming Interface}
  \item[BLE] \textit{Bluetooth Low Energy}
  \item[CNN] \textit{Convolutional Neural Network}
  \item[DSP] \textit{Digital Signal Processing} (Processamento Digital de Sinais)
  \item[FIFO] \textit{First-In, First-Out}
  \item[HAR] \textit{Human Activity Recognition} (Reconhecimento de Atividades Humanas)
  \item[IMU] \textit{Inertial Measurement Unit} (Unidade de Medição Inercial)
  \item[LOSO] \textit{Leave-One-Subject-Out}
  \item[LSTM] \textit{Long Short-Term Memory}
  \item[MAE] \textit{Mean Absolute Error} (Erro Médio Absoluto)
  \item[ML] \textit{Machine Learning} (Aprendizado de Máquina)
  \item[RMSE] \textit{Root Mean Squared Error} (Raiz do Erro Quadrático Médio)
  \item[SDK] \textit{Software Development Kit}
  \item[SVM] \textit{Support Vector Machine}
  \item[TCC] Trabalho de Conclusão de Curso
  \item[UI] \textit{User Interface} (Interface de Usuário)
\end{siglas}

% Lista de simbolos
\begin{simbolos}
  \item[$\mathbf{a}(t)$] Vetor de aceleração inercial no instante $t$
  \item[$\boldsymbol{\omega}(t)$] Vetor de velocidade angular (giroscópio) no instante $t$
  \item[$\|\mathbf{a}(t)\|$] Magnitude da aceleração inercial
  \item[$\|\boldsymbol{\omega}(t)\|$] Magnitude da velocidade angular
  \item[$f_s$] Frequência de amostragem ($f_s = 25\text{ Hz}$, $T_s = 0{,}04\text{ s}$)
  \item[$f_c$] Frequência de corte do filtro passa-baixa ($f_c \in [3, 5]\text{ Hz}$)
  \item[$t_{min}$] Intervalo temporal mínimo entre repetições válidas ($t_{min} \approx 1{,}2\text{ s}$)
\end{simbolos}

% Sumario
\pdfbookmark[0]{\contentsname}{toc}
\tableofcontents*
\cleardoublepage

% ELEMENTOS TEXTUAIS
\textual

% CAPITULO 1: INTRODUCAO
\chapter{Introdução}

\section{Contextualização e Motivação}
A prática regular de exercícios resistidos (musculação) é amplamente recomendada por diretrizes de saúde internacionais devido aos seus benefícios comprovados no ganho de massa muscular, aumento da densidade óssea e prevenção de doenças metabólicas. Para maximizar os resultados e garantir a segurança dos praticantes, o controle de variáveis do treinamento - como número de séries, repetições executadas, cadência e tempo sob tensão - é indispensável.

Com a massificação de dispositivos vestíveis (\textit{smartwatches}), como o Apple Watch, sensores inerciais de alta precisão (acelerômetros e giroscópios) tornaram-se onipresentes no pulso dos usuários. Todavia, a maioria dos aplicativos de saúde comercialmente disponíveis limita-se ao registro de parâmetros cardiovasculares gerais e estimativas calóricas, negligenciando a contagem automática e fidedigna de repetições durante o treino resistido.

\section{Formulação do Problema}
A identificação precisa de repetições em tempo real a partir de sensores de pulso impõe desafios computacionais e biomecânicos significativos:
\begin{enumerate}
    \item \textbf{Variabilidade Intra e Inter-Sujeito:} Diferentes usuários executam o mesmo exercício com cadências, amplitudes articulares e biotipos distintos.
    \item \textbf{Ruídos Inerciais e Movimentos Parasitas:} Movimentos involuntários, tremores sob fadiga e ajustes de pegada geram perturbações no sinal que podem ser erroneamente interpretadas como repetições.
    \item \textbf{Restrições de Recursos:} O processamento em tempo real deve operar com consumo energético moderado e latência imperceptível ao usuário, inviabilizando abordagens computacionalmente proibitivas em dispositivos móveis.
\end{enumerate}

Dessa forma, formula-se a questão de pesquisa: \textit{Como estruturar um sistema móvel e vestível capaz de processar dados inerciais em tempo real para contar repetições de exercícios resistidos com alta precisão e baixa latência?}

\section{Objetivos}

\subsection{Objetivo Geral}
Desenvolver e validar uma arquitetura de \textit{software} móvel e vestível (Gymbeat) capaz de coletar sinais inerciais a partir de um \textit{smartwatch}, transmiti-los ao \textit{smartphone} e executar um modelo híbrido de Inteligência Artificial e Processamento Digital de Sinais para contagem de repetições em tempo real durante exercícios de musculação.

\subsection{Objetivos Específicos}
\begin{enumerate}
    \item Implementar um módulo no watchOS configurado a 25 Hz para aquisição contínua de aceleração e velocidade angular.
    \item Construir um protocolo de sincronização local e em nuvem via Firebase Firestore para consolidação de um \textit{dataset} de treino anotado com \textit{ground truth}.
    \item Projetar e implementar um pipeline de Processamento Digital de Sinais baseado em filtragem passa-baixa Butterworth e magnitude vetorial.
    \item Desenvolver um algoritmo de detecção de extremos locais e validação por Aprendizado de Máquina clássico para contagem de repetições.
    \item Validar o algoritmo através do protocolo experimental \textit{Leave-One-Subject-Out} (LOSO), mensurando $MAE$, $RMSE$ e latência de inferência no iPhone.
\end{enumerate}

\section{Justificativa e Relevância}
Este trabalho contribui para o estado da arte em Reconhecimento de Atividades Humanas (HAR) ao propor uma solução leve, escalável e de fácil adoção prática. Ao transferir a carga computacional para o smartphone e utilizar um modelo híbrido eficiente, o sistema supera as limitações de bateria dos vestíveis e elimina a necessidade de \textit{hardware} adicional dedicado em academias.

% CAPITULO 2: FUNDAMENTACAO TEORICA
\chapter{Fundamentação Teórica e Trabalhos Relacionados}

\section{Computação Vestível e Unidades de Medição Inercial (IMU)}
Uma Unidade de Medição Inercial (IMU) combina acelerômetros triaxiais e giroscópios triaxiais:
\begin{itemize}
    \item \textbf{Acelerômetro Triaxial:} Mede a aceleração própria somada à gravidade nas componentes ortogonais $a_x(t), a_y(t), a_z(t)$ expressas em $g$ ($1g \approx 9{,}81\text{ m/s}^2$).
    \item \textbf{Giroscópio Triaxial:} Mensura a velocidade angular $\omega_x(t), \omega_y(t), \omega_z(t)$ expressa em radianos por segundo ($\text{rad/s}$) ou graus por segundo ($\text{°/s}$).
\end{itemize}

\section{Biomecânica dos Exercícios Resistidos e Fases do Movimento}
Os exercícios de musculação são caracterizados por ciclos periódicos compostos tipicamente por:
\begin{enumerate}
    \item \textbf{Fase Excêntrica:} O músculo-alvo se alonga sob tensão (ex: descida da barra no supino ou agachamento).
    \item \textbf{Ponto de Inflexão (Transição):} Instante de velocidade linear mínima onde ocorre a inversão do movimento.
    \item \textbf{Fase Concêntrica:} O músculo-alvo se contrai superando a resistência (ex: subida da carga).
\end{enumerate}

\section{Processamento Digital de Sinais Aplicado a Sensores}

\subsection{Invariância de Orientação via Magnitude Vetorial}
Para garantir invariância espacial em relação a pequenas rotações do relógio no pulso, calcula-se a magnitude euclidiana:
\begin{equation}
    \|\mathbf{a}(t)\| = \sqrt{a_x(t)^2 + a_y(t)^2 + a_z(t)^2}
\end{equation}
\begin{equation}
    \|\boldsymbol{\omega}(t)\| = \sqrt{\omega_x(t)^2 + \omega_y(t)^2 + \omega_z(t)^2}
\end{equation}

\subsection{Filtragem Passa-Baixa Butterworth}
O movimento humano voluntário concentra-se abaixo de 3 a 5 Hz. Ruídos de alta frequência decorrentes de tremores musculares ou imperfeições eletrônicas são atenuados por um filtro digital Butterworth de 4ª ordem.

\section{Aprendizado de Máquina em Reconhecimento de Atividades (HAR)}
Modelos clássicos como \textit{Random Forests} e \textit{Support Vector Machines} (SVM) oferecem alta capacidade discriminativa com baixo custo computacional para vetores de atributos extraídos de picos inerciais.

% CAPITULO 3: METODOLOGIA E ARQUITETURA
\chapter{Metodologia e Arquitetura Proposta}

\section{Arquitetura Geral do Sistema (Gymbeat)}
O sistema opera em três camadas: Apple Watch (coletor a 25 Hz), iPhone (motor de inferência, buffer circular e interface gráfica) e Firebase Firestore (armazenamento e histórico em nuvem).

\section{Decisão Arquitetural: Descarte do Deep Learning Fim-a-Fim}
A abordagem de Deep Learning fim-a-fim (1D-CNN / BiLSTM) foi descontinuada devido à ausência de anotação frame-a-frame (o protocolo coleta repetições totais por série) e à necessidade de datasets massivos para evitar \textit{overfitting}. Adotou-se a abordagem híbrida de Processamento de Sinais + Validação por IA Clássica.

\section{Algoritmo Híbrido Proposto}
O algoritmo opera sobre uma janela deslizante em buffer circular, detectando inflexões de aceleração e validando cada evento por meio do classificador treinado.

\begin{algorithm}[H]
\caption{Contagem e Validação em Tempo Real (Gymbeat)}
\begin{algorithmic}[1]
\REQUIRE Fluxo contínuo de amostras $s_t = (a_x, a_y, a_z, \omega_x, \omega_y, \omega_z)$
\ENSURE $reps$ (contagem acumulada de repetições válidas)
\STATE Inicializar buffer circular $B$ com capacidade $W = 75$ amostras
\STATE $reps \leftarrow 0$, $t_{last} \leftarrow 0$, $t_{min} \leftarrow 1{,}2\text{ s}$
\FOR{cada amostra $s_t$ recebida}
    \STATE $mag \leftarrow \sqrt{a_x^2 + a_y^2 + a_z^2}$
    \STATE $mag_{filt} \leftarrow \text{ButterworthPassaBaixa}(mag)$
    \STATE $B.\text{adicionar}([s_t.timestamp, mag_{filt}, \boldsymbol{\omega}])$
    \IF{$B.\text{isFull}()$}
        \IF{$\text{DetectarExtremoLocal}(B)$ \AND $(s_t.timestamp - t_{last} \ge t_{min})$}
            \STATE $features \leftarrow \text{ExtrairAtributosPico}(B)$
            \STATE $classe \leftarrow \text{ClassificadorIA}(features)$
            \IF{$classe == \text{VALIDO}$}
                \STATE $reps \leftarrow reps + 1$
                \STATE $t_{last} \leftarrow s_t.timestamp$
                \STATE $\text{EmitirFeedback}(reps)$
            \ENDIF
        \ENDIF
    \ENDIF
\ENDFOR
\end{algorithmic}
\end{algorithm}

\section{Métricas de Avaliação}
\begin{equation}
    \text{MAE} = \frac{1}{M} \sum_{i=1}^M |\text{reps\_real}_i - \text{reps\_pred}_i|
\end{equation}
\begin{equation}
    \text{RMSE} = \sqrt{\frac{1}{M} \sum_{i=1}^M (\text{reps\_real}_i - \text{reps\_pred}_i)^2}
\end{equation}

% CAPITULO 4: DESENVOLVIMENTO
\chapter{Desenvolvimento e Implementação}

\section{Coletor no Apple Watch (Swift)}
Configurado a 25 Hz ($T_s = 0{,}04\text{ s}$) via \texttt{CMMotionManager} e \texttt{HealthKit} no arquivo \texttt{WorkoutManager.swift}.

\section{Transmissão e Sincronização}
Transmissão sem fio via \texttt{WCSession} (\texttt{WatchConnectivity}) e persistência em lote no Firebase Firestore.

\section{Motor de Inferência no iPhone}
Implementado em TypeScript no módulo \texttt{src/ai/RepetitionCounter.ts} utilizando buffer circular FIFO (\texttt{CircularBuffer.ts}).

% CAPITULO 5: RESULTADOS E DISCUSSAO
\chapter{Resultados e Discussões}

\section{Avaliação Experimental Preliminar}
Resultados preliminares indicam latência de processamento inferior a 5 ms no dispositivo móvel, com taxa de amostragem de 25 Hz suficiente para captura dos extremos biomecânicos.

% CAPITULO 6: CONCLUSAO
\chapter{Conclusão e Trabalhos Futuros}

\section{Considerações Finais}
A arquitetura híbrida desenvolvida demonstrou alto potencial de aplicação prática, unindo facilidade de uso do smartwatch com a robustez de processamento do smartphone.

% ELEMENTOS POS-TEXTUAIS
\postextual

% Referencias bibliograficas
\begin{thebibliography}{99}

\bibitem{batchelor2020}
BATCHELOR, R. et al. Automatic exercise recognition and repetition counting using smartwatches. \textbf{IEEE Journal of Biomedical and Health Informatics}, v. 24, n. 3, p. 892-901, 2020.

\bibitem{bulling2014}
BULLING, A.; BLANKE, U.; SCHIELE, B. A tutorial on human activity recognition using body-worn inertial sensors. \textbf{ACM Computing Surveys (CSUR)}, v. 46, n. 3, p. 1-33, 2014.

\bibitem{lara2013}
LARA, O. D.; LABRADOR, M. A. A survey on human activity recognition using wearable sensors. \textbf{IEEE Communications Surveys \& Tutorials}, v. 15, n. 3, p. 1192-1209, 2013.

\bibitem{oppenheim2009}
OPPENHEIM, A. V.; SCHAFER, R. W. \textbf{Discrete-Time Signal Processing}. 3. ed. Upper Saddle River: Pearson, 2009.

\bibitem{pedregosa2011}
PEDREGOSA, F. et al. Scikit-learn: Machine learning in Python. \textbf{Journal of Machine Learning Research}, v. 12, p. 2825-2830, 2011.

\bibitem{shoeb2021}
SHOEB, M. et al. Real-time repetition counting and form analysis in strength training using IMU sensors. In: \textbf{Proceedings of the ACM on Interactive, Mobile, Wearable and Ubiquitous Technologies}, v. 5, n. 2, p. 1-22, 2021.

\end{thebibliography}

\end{document}
